\section{结论}

本文的主线是把有限窗口读出所诱导的微态空间，通过可审计的构造链条组织为稳定类型、可实现的算术结构与可证书化的统计/谱指纹接口。相应的主要结论可按“输入口径—构造机制—桥接与验证”三段收束如下。

\subsection{主要结论：按机制分组}\label{subsec:core-conclusions}
\paragraph{输入口径的固定}
数论侧以反共振规范代表与有限样本差异度证书固定稳定性基线（第 \ref{sec:principles} 节），动力学侧以扫描—投影生成器（群作用与二元读出）固定可观测对象的最小输入（第 \ref{sec:spg} 节）。上述口径在全文保持不变，使后续结论具有可追溯的依赖链。
在该接口上，“扫描深度 \(\Rightarrow\) 清晰度”可被压缩为最优判错剖面与边界指数律，并在地址语义下进一步分解为“未同步（统计型 \(\mathsf{NULL}\)）”与“地址碰撞”两类失效机制的双预算容量律（定理 \ref{thm:spg-clarity-boundary-dimension}、\ref{thm:spg-double-budget-address-capacity}）；在在线制度下，该清晰度剖面又与同步扣除的 Chebotarev 容量不等式耦合，从而把定义性同步门槛与统计均匀化预算分离为可审计的两段账本（定理 \ref{thm:pom-sync-subtracted-chebotarev-capacity}）。

\paragraph{分辨率折叠与递归地址化}
本文构造有限窗口折叠 \(\Fold_m\) 并证明其跨分辨率相容与统一的 sofic 因子接口结论（定理 \ref{thm:fold-suite}；并见定义 \ref{def:fold-word} 与命题 \ref{prop:fold-basic}）。在同一可测框架内，递归地址化被严格化为“可读出前缀”的内生化机制：地址对象未给定时，评价函数 \(\ind_{C_a}\) 在指称层不成立，读出偏函数不定义并以 \(\mathsf{NULL}\) 表示不可寻址；一旦地址给定，派生柱事件族不扩张既有可见 \(\sigma\)-代数，而仅在旧层可见 \(\mathcal{G}^{(L)}\) 内重选并组织可见事件，从而 \(\mathcal{G}^{(L+1)}\subseteq \mathcal{G}^{(L)}\)（命题 \ref{prop:recursive-addressing-refinement}；推论 \ref{cor:recursive-addressing-visible-sigma-nonexpansion}）。折叠降阶的非局域性可由规范差 \(G_m\) 定量审计并给出最坏与典型的双重锚点（定义 \ref{def:fold-gauge-anomaly}；定理 \ref{thm:fold-gauge-anomaly-max}、\ref{thm:fold-gauge-anomaly-density}）。

\paragraph{稳定算术与投影本体数学}
在值保持重写与规范横截面机制下，分辨率 \(m\) 的稳定类型算术与 \(\ZZ/(F_{m+2}\ZZ)\) 同构（定理 \ref{thm:finite-resolution-mod}），从而把统一/分裂的第一性来源归结为模数的环论性质及 CRT 结构（推论 \ref{cor:field-phase-fib-prime}、\ref{cor:crt-factorization}）；并可沿二幂子塔定义 Fibonacci-adic 逆极限环并得到紧致单位群（定义 \ref{def:fib-adic-ring}；命题 \ref{prop:fib-adic-basic}）。在此基础上，本文建立投影本体数学（POM）的投影词演算，并给出信息账本恒等式，将粗粒化的信息损失严格对象化为可计算残差（定理 \ref{thm:pom-kl-ledger}；推论 \ref{cor:pom-ent-increase-equals-kl}）。同时在局部读出塔片段上证明重写的终止性与合流性并给出正规形，进而得到可审计闭环与判定性接口（定理 \ref{thm:pom-rewrite-termination-confluence}；命题 \ref{prop:pom-projword-decidable}、\ref{prop:pom-rewrite-audit-loop}）。高阶碰撞矩的 moment-kernel 编译与共振判别把“共振”落实为可检验机制（定理 \ref{thm:pom-collision-kernel-ak}、\ref{thm:pom-moment-resonance}）。

\paragraph{桥接接口与验证}
在谱侧，动力学 \(\zeta\) 与有限部接口把周期计数、迹与主谱半径组织为可审计指纹，并为后续的完成化与扭曲通道提供统一输入（第 \ref{sec:zeta-finite-part} 节）。在统计侧，折叠直方图的偏差指标、跨分辨率一致性残差与到 Parry(PF) 基线的误差分解式将有限样本误差与模型差异项显式分离（第 \ref{sec:statistical-stability}、\ref{sec:experiments} 节；推论 \ref{cor:tv-decomposition-parry}），并以脚本化协议给出可复现实证书输出（第 \ref{sec:experiments} 节）。
其中在 \(\zeta/\Xi\) 的完成化审计中，离开 unitary slice 的信息被端点二次径向压缩并导出严格的可见化预算门槛与单半径带状禁入证书，同时最小 \((2\times2)\) 自对偶模板给出显式轴外零点塔作为可证伪原型（定理 \ref{thm:xi-jensen-single-radius-band-exclusion}、\ref{thm:xi-2x2-selfdual-offcritical-template}）。

\paragraph{四项结论性命题：从地址语义到证书化预算与秩稳定性}
为便于把正文中分散出现的结论直接复用到协议设计与审计流程中，可将若干核心断言压缩为四条互补的“证书接口”（均严格依赖既有编号结论而非解释性表述）：
\begin{enumerate}
  \item \textbf{地址语义的对象化：\(\mathsf{NULL}\) 作为局部截面障碍。}
  在前缀站点上把协议可接受且单值可核验的读出/证书组织为（预）层 \(\mathscr F\) 后，“地址先于取值”的 \(\mathsf{NULL}\) 语义可被严格等价为 \(\mathscr F(a)=\varnothing\) 的存在性障碍，并与递归地址化的可见 \(\sigma\)-代数不扩张口径相容（\S\ref{subsec:prefix-site-null}，命题 \ref{prop:null-as-local-section-obstruction}；并见命题 \ref{prop:recursive-addressing-refinement} 与推论 \ref{cor:recursive-addressing-visible-sigma-nonexpansion}）。
  由此，语义/协议/碰撞三类空值在同一站点语言下被统一为三类互不混淆的局部截面缺失机制（推论 \ref{cor:null-trichotomy-local-section}）。

  \item \textbf{折叠“清晰度”的二阶证书：碰撞膨胀 \(\Rightarrow\) 熵缺口上界。}
  在均匀微态基线下，折叠推前分布 \(w_m(x)=d_m(x)/2^m\) 相对类型均匀分布的熵缺口 \(\mathrm{gap}\) 可被二阶碰撞指纹
  \(C_{\mathrm{col}}(m)=|X_m|\sum_x w_m(x)^2\)
  的对数上界证书化（定理 \ref{thm:pom-spectral-information-sandwich} 与定义 \ref{def:pom-collision-inflation}），并可由表 \ref{tab:fold_multiplicity_histogram} 在有限分辨率上直接审计。
  该结论将“缩略图的清晰度”从解释性说法压缩为一个仅依赖二阶统计量的可计算证书。

  \item \textbf{单射化预算的最优对偶包络：分位配额律的 \(\inf_{q\ge 2}\) 规范化。}
  prime-register 的配额阈值可由纤维指数尾部分位 \(B_m(\varepsilon)=\exp(m\gamma_m(\varepsilon))\) 给出，并在区间读出/重建口径下具有锐性（\S\ref{subsec:pom-quantile-budget-law}，推论 \ref{cor:pom-prime-register-quantile-budget}）。
  进一步，整数矩阶 \(q\) 所导出的全部 Chernoff/配额不等式族可被压缩为对 \(q\ge 2\) 的单一最优下包络，从而得到一条参数自由的寄存器指数预算律；并可用有限 \(q\)-窗口的“最强可计算尾界”实现可执行证书（推论 \ref{cor:pom-prime-register-q-quota-family} 及 \S\ref{subsec:pom-quantile-budget-law} 的对偶设计式）。

  \item \textbf{模式轴完备性的充要判据：Hankel 秩一致有界 \(\Longleftrightarrow\) ACC-Rank。}
  在统一最小性与协议闭包假设下，一旦残差轴预算已排除碰撞空值，则“不会因需要新增可见自由度而触发 \(\mathsf{NULL}\)”的模式轴完备性与 Hankel 秩的一致有界性严格等价（定理 \ref{thm:xi-completeness-iff-hankel-bounded}），从而把 completeness 的关键缺口转换为可回归测试的秩稳定性扫描（注 \ref{rem:xi-acc-rank}）。
  与分位配额证书合并后，completeness 的审计任务可被压缩为“秩稳定证书 \(\times\) 配额证书”的二元后端核验（命题 \ref{prop:xi-completeness-auditable-decomposition}）。
\end{enumerate}

\subsection{复现与审计指引}\label{subsec:repro-audit-pointer}
解释性综合、边界与可证伪点集中于第 \ref{sec:discussion} 节。附录汇集算子代数、相位塔、同步核与多尺度接口等支撑材料，以支撑正文的可审计推理链条与复核需求。
